% !TeX program = xelatex
% OVERLEAF: Settings -> Compiler -> XeLaTeX. Dòng !TeX không tự đổi compiler.
% Mẫu slide tiến độ gần một tháng; chỉ đưa hình/bảng từ kết quả thực.
\documentclass[10pt,aspectratio=169]{beamer}
\usetheme{Madrid}
\usecolortheme{seahorse}
\usepackage{fontspec}
\IfFontExistsTF{DejaVu Sans}{\setsansfont{DejaVu Sans}}{\setsansfont{Arial}}
\IfFontExistsTF{DejaVu Serif}{\setmainfont{DejaVu Serif}}{\setmainfont{Times New Roman}}
\usepackage{amsmath,amssymb,graphicx,booktabs,xcolor}
\setbeamertemplate{navigation symbols}{}
\setbeamersize{text margin left=7mm,text margin right=7mm}
\newcommand{\todo}[1]{{\color{red!70!black}[Điền: #1]}}
\newcommand{\imagebox}[1]{\fbox{\parbox[c][3.1cm][c]{.92\linewidth}{\centering #1\\\small Thay bằng hình kết quả thực}}}
\title{Từ FEM đến PINN dạng mạnh và mixed PINN cho bài toán dầm}
\subtitle{Tiến độ nghiên cứu sau gần một tháng}
\author{Nguyễn Lin Đa \quad Trà Minh Thiện \quad Võ Trần Tiến}
\institute{Hướng dẫn: TS. Nguyễn Minh Nhân}
\date{Tháng 11 năm 2026}
\begin{document}

\begin{frame}
  \titlepage
  \vspace{-4mm}
  \begin{center}\small Chỉ trình bày các phần nhóm đã thực sự hoàn thành.\end{center}
\end{frame}

\begin{frame}{1. Mục tiêu và lộ trình nghiên cứu}
\begin{block}{Câu hỏi}
PINN dạng mạnh giải bài toán dầm chính xác và tốn chi phí ra sao so với FEM? Dạng mixed có cải thiện được các điểm yếu quan sát thấy không?
\end{block}
\begin{enumerate}
  \item Học cơ sở PINN; dựng nghiệm giải tích và FEM Hermite làm mốc.
  \item Cài PINN dạng mạnh; đo ưu, nhược điểm bằng kết quả.
  \item Thử mixed PINN trên cùng bài toán và đối chiếu công bằng.
\end{enumerate}
\small\todo{Đánh dấu bước đã hoàn thành và bước còn đang thực hiện.}
\end{frame}

\begin{frame}{2. Tài liệu nền và câu hỏi cần kiểm chứng}
\small
\begin{tabular}{p{.18\linewidth}p{.37\linewidth}p{.37\linewidth}}
\toprule
Phương pháp & Nguồn đọc & Vai trò trong nghiên cứu\\
\midrule
FEM Hermite & Reddy (2019), phần tử dầm Euler--Bernoulli & Dạng yếu, hàm Hermite và hội tụ theo lưới.\\
PINN mạnh & Raissi et al. (2019); Katsikis et al. (2022) & Residual phương trình và điều kiện biên; bài toán dầm tĩnh.\\
Mixed PINN & Mouratidou et al. (2026) & Ý tưởng dùng trường phụ cho dầm \emph{động}; nhóm thử cho dầm \emph{tĩnh}.\\
\bottomrule
\end{tabular}
\vspace{3mm}
\begin{block}{Câu hỏi}
Giảm bậc đạo hàm trong mixed PINN có thực sự cải thiện sai số mômen, thời gian hoặc độ ổn định so với PINN mạnh trên bài toán này không?
\end{block}
\end{frame}

\begin{frame}{3. Bài toán và nghiệm tham chiếu}
\begin{columns}[T,onlytextwidth]
\begin{column}{.54\textwidth}
\begin{align*}
 (D(x)w'')''&=q_0,\\
 w(0)=w(L)&=0,\\
 m(0)=m(L)&=0,\qquad m=Dw''.
\end{align*}
Với $L=q_0=D=1$:
\[w_*(x)=\frac{x^4-2x^3+x}{24},\qquad m_*(x)=\frac{x^2-x}{2}.\]
\end{column}
\begin{column}{.43\textwidth}
\begin{alertblock}{Điều cần kiểm tra}
Nghiệm tham chiếu thỏa phương trình và bốn điều kiện biên. Nghiệm này chỉ dùng để đánh giá sau khi giải.
\end{alertblock}
\small\todo{Nếu làm dầm hai đoạn, nêu $D(x)$ và bốn điều kiện ghép.}
\end{column}
\end{columns}
\end{frame}

\begin{frame}{4. FEM Hermite làm mốc đối chiếu}
\begin{columns}[T,onlytextwidth]
\begin{column}{.56\textwidth}
\begin{block}{Từ dạng yếu đến hệ đại số}
\[\int_0^L Dw_h''v_h''\,dx=\int_0^L q_0v_h\,dx,\qquad
 \mathbf K\mathbf d=\mathbf f.\]
\[\mathbf K_e=\int_e D\mathbf B^{\mathsf T}\mathbf B\,dx,\qquad
 \mathbf f_e=\int_e q_0\mathbf N^{\mathsf T}\,dx.\]
\end{block}
\small Hai bậc tự do mỗi nút: $w_i$ và góc xoay $\theta_i$. Cố định $w$ ở hai đầu, để góc xoay tự do.
\end{column}
\begin{column}{.41\textwidth}
\begin{block}{Kiểm tra bắt buộc}
\begin{itemize}
\item Lưới $n_e=\todo{...}$.
\item So với $w_*,m_*$.
\item Ít nhất ba mức lưới để xem hội tụ.
\end{itemize}
\end{block}
\end{column}
\end{columns}
\end{frame}

\begin{frame}{5. PINN dạng mạnh: baseline cần đánh giá trước}
\begin{columns}[T,onlytextwidth]
\begin{column}{.55\textwidth}
\[x\longmapsto w_\theta(x),\qquad
  r_{\rm strong}(x)=(Dw_\theta'')''-q_0.\]
\[\mathcal L_{\rm strong}=\lambda_r\operatorname{MSE}(r_{\rm strong})
 +\lambda_b\!\sum_{x\in\{0,L\}}\!\bigl[w_\theta(x)^2+(Dw_\theta'')(x)^2\bigr].\]
\small Với $D$ không đổi, residual cần đạo hàm bậc bốn của $w_\theta$.
\end{column}
\begin{column}{.42\textwidth}
\begin{block}{Cần báo cáo}
\begin{itemize}
\item Sai số $w$ và $m=Dw_\theta''$.
\item Sai lệch biên, residual trên lưới độc lập.
\item Thời gian train, số bước và độ ổn định theo seed.
\end{itemize}
\end{block}
\end{column}
\end{columns}
\end{frame}

\begin{frame}{6. Mixed PINN: thử nghiệm sau baseline}
\begin{columns}[T,onlytextwidth]
\begin{column}{.54\textwidth}
\[x\longmapsto\bigl(w_\theta(x),m_\theta(x)\bigr)\]
\[r_c=m_\theta-Dw_\theta'',\qquad r_e=m_\theta''-q_0.\]
\[\mathcal L=\lambda_c\operatorname{MSE}(r_c)
 +\lambda_e\operatorname{MSE}(r_e)+\lambda_b\mathcal L_{\rm biên}.\]
\small Residual chỉ cần đạo hàm bậc hai. Cần kiểm tra bằng số liệu xem điều này có giúp train và cải thiện $w,m$ hay không.
\end{column}
\begin{column}{.43\textwidth}
\begin{block}{Cấu hình đã dùng}
\begin{itemize}
\item Mạng/activation: \todo{...}
\item Điểm collocation: \todo{...}
\item Optimizer, bước train: \todo{...}
\item Trọng số loss và seed: \todo{...}
\end{itemize}
\end{block}
\end{column}
\end{columns}
\end{frame}

\begin{frame}{7. Thiết kế so sánh và thước đo sai số}
\small
\begin{columns}[T,onlytextwidth]
\begin{column}{.51\textwidth}
\begin{block}{Đánh giá trên lưới độc lập}
Với $g=w$ hoặc $m$ và nghiệm phương pháp $g_j$:
\[E_{L^2}^{(j)}(g)=
\left(\frac{\int_0^L|g_j-g_*|^2\,dx}
{\int_0^L|g_*|^2\,dx}\right)^{1/2}.\]
\[E_{\infty}^{(j)}(g)=
\frac{\max_{X_{\rm test}}|g_j-g_*|}
{\max_{X_{\rm test}}|g_*|}.\]
\end{block}
\end{column}
\begin{column}{.46\textwidth}
\begin{block}{Cùng điều kiện thí nghiệm}
\begin{itemize}
\item Cùng bài toán và lưới kiểm tra.
\item Sai lệch biên $w,m$; sai số ở giữa dầm.
\item FEM: hội tụ theo $n_e$; PINN: nhiều seed.
\item Thời gian lắp ráp/giải hoặc train/dự đoán.
\end{itemize}
\end{block}
\end{column}
\end{columns}
\small Loss train và residual được báo cáo riêng, không xem là sai số nghiệm.
\end{frame}

\begin{frame}{8. Kết quả 1: đường võng và mômen}
\begin{columns}[T,onlytextwidth]
\begin{column}{.49\textwidth}\centering\imagebox{Đồ thị $w_*,w_{\rm FEM},w_{\rm strong},w_{\rm mixed}$}\end{column}
\begin{column}{.49\textwidth}\centering\imagebox{Đồ thị $m_*,m_{\rm FEM},m_{\rm strong},m_{\rm mixed}$}\end{column}
\end{columns}
\vspace{2mm}\small\todo{Bỏ đường mixed nếu chưa chạy; ghi cấu hình, seed và nơi sai khác lớn nhất.}
\end{frame}

\begin{frame}{9. Kết quả 2: phân bố sai số và hội tụ}
\begin{columns}[T,onlytextwidth]
\begin{column}{.49\textwidth}\centering\imagebox{Đồ thị $|w_j-w_*|$ và $|m_j-m_*|$ theo $x$}\end{column}
\begin{column}{.49\textwidth}\centering\imagebox{FEM: sai số theo $n_e$; PINN: loss và sai số test theo bước train}\end{column}
\end{columns}
\vspace{2mm}\small\todo{Nêu nơi sai số lớn nhất; phân biệt loss train với sai số test.}
\end{frame}

\begin{frame}{10. Kết quả 3: sai số và chi phí}
\small
\begin{table}
\centering
\begin{tabular}{lccccc}
\toprule
Phương pháp & $E_{L^2}(w)$ & $E_{L^2}(m)$ & $E_\infty(w)$ & Lệch biên & Thời gian (s)\\
\midrule
FEM Hermite & --- & --- & --- & --- & ---\\
PINN dạng mạnh & --- & --- & --- & --- & ---\\
Mixed PINN (nếu đã chạy) & --- & --- & --- & --- & ---\\
\bottomrule
\end{tabular}
\end{table}
\begin{itemize}
\item Lưới kiểm tra: \todo{số điểm}; phần cứng/dtype: \todo{...}; PINN: \todo{số seed}.
\item Tách thời gian FEM lắp ráp/giải và PINN train/dự đoán trong báo cáo chi tiết.
\item So sánh tại mức sai số đạt được; bỏ hàng của phương pháp chưa chạy.
\end{itemize}
\end{frame}

\begin{frame}{11. Kết luận có giới hạn và bước tiếp theo}
\begin{block}{Nhóm đã chứng minh được gì bằng số liệu?}
\todo{Nêu ưu, nhược điểm của PINN dạng mạnh bằng sai số, residual và thời gian; nếu đã chạy, cho biết mixed thay đổi gì.}
\end{block}
\begin{alertblock}{Chưa thể kết luận gì?}
\todo{Ví dụ: mới một seed, mixed chưa chạy đủ cấu hình, hoặc chưa kiểm tra độ cứng gián đoạn.}
\end{alertblock}
\begin{itemize}
\item Trước mắt: hoàn thành so sánh strong PINN với FEM rồi thử mixed PINN cùng cấu hình.
\item Sau đó: nhiều seed; nếu còn thời gian, $D(x)$ hai đoạn và điều kiện ghép tại $a$.
\end{itemize}
\end{frame}

\begin{frame}{Tài liệu và mã để tái lập}
\small
\begin{enumerate}
\item Reddy (2019), \emph{Introduction to the Finite Element Method}, 4th ed., phần tử dầm Euler--Bernoulli.
\item Raissi, Perdikaris, Karniadakis (2019), \emph{Physics-informed neural networks}, DOI: \href{https://doi.org/10.1016/j.jcp.2018.10.045}{10.1016/j.jcp.2018.10.045}.
\item Katsikis, Muradova, Stavroulakis (2022), \emph{A Gentle Introduction to Physics-Informed Neural Networks...}, DOI: \href{https://doi.org/10.15377/2409-5761.2022.09.8}{10.15377/2409-5761.2022.09.8}.
\item Mouratidou et al. (2026), \emph{A Mixed Formulation Physics-Informed Neural Network...}, DOI: \href{https://doi.org/10.3390/a19100851}{10.3390/a19100851}.
\item Mã, seed và dữ liệu kết quả: \todo{đường dẫn repository/thư mục của nhóm}.
\end{enumerate}
\vspace{5mm}
\begin{center}\Large Cảm ơn thầy và các bạn!\end{center}
\end{frame}
\end{document}
